% The class loads wasysym with its default `compat1` option, which \lets
% \iint/\iiint to wasysym's own symbols at package-load time (i.e. before
% mhchem's internal amsmath load happens later in the same class). amsmath
% then tries \newcommand{\iint} and fails because the name is already taken
% -- a real clash between this class's dependencies on current TeX Live, not
% something introduced here. `nointegrals` keeps wasysym from touching
% \iint/\iiint at all, which is the settings a text-only newsletter needs.
\PassOptionsToPackage{nointegrals}{wasysym}

\documentclass{astropah}

\graphicspath{{images}}

% Information that needs to be updated in every issue
\newcommand{\monthyear}{September 2026}
\newcommand{\issuenum}{1}
\newcommand{\nextissue}{15 October 2026}
\newcommand{\nextdeadline}{1 October 2026}

\begin{document}

% =================================================
%  EDITORIAL
% =================================================
\begin{editorial}

  Welcome to the first issue of this newsletter. Polycyclic Aromatic
  Hydrocarbons (PAHs) remain one of the most abundant classes of complex
  organic molecules detected in space, carrying an estimated 10--20\% of the
  carbon in the interstellar medium. This issue rounds up three recent
  papers on PAH emission features, a JWST spectroscopy campaign, and the
  usual meetings calendar.

  As always, we welcome short summaries of newly accepted papers, PhD
  positions, and meeting announcements from the community for the next
  issue.

\end{editorial}

\vfill

\usefont {OT1} {phv}{b}{n}
\fontsize{15}{15}
\selectfont

\begin{shaded}
  \vspace{0.1cm}
  \centerline{\textcolor{white}{Next issue: \textbf{\nextissue}.}}
  \vspace{0.2cm}
  \centerline{\textcolor{white}{Submission deadline: \textbf{\nextdeadline}.}}
  \vspace{0.1cm}
\end{shaded}

\newpage

% =================================================
%  In Focus
% =================================================
\begin{infocus}

  \begin{center}
    {\infocustitle \noindent \purple{Aromatic Infrared Bands Toward the
    Orion Bar}}
  \end{center}
  \vspace{0.4cm}
  \normalsize

  The Orion Bar is a favourite laboratory for PAH astrochemistry: a dense
  photodissociation region illuminated almost edge-on, so the transition
  from ionised gas through the PAH-emission layer to the molecular cloud
  can be mapped spatially rather than inferred from an unresolved spectrum.

  Recent JWST/NIRSpec integral-field observations resolve the 3.3, 6.2 and
  7.7 \textmu m aromatic infrared bands across this transition at
  sub-arcsecond resolution for the first time. The band ratios shift
  systematically with distance from the ionisation front, consistent with a
  changing balance between PAH cations and neutrals as the local radiation
  field softens.

  \begin{itemize}
    \item The 3.3\,\textmu m/11.2\,\textmu m ratio, a common proxy for PAH
      ionisation state, drops by a factor of $\sim$3 across the mapped
      region.
    \item No strong evidence yet for a spatial break in PAH size
      distribution, contrary to some pre-JWST predictions.
    \item Follow-up MIRI/MRS spectroscopy is scheduled to extend the map to
      the 12--18\,\textmu m range this cycle.
  \end{itemize}

\end{infocus}

% =================================================
%  Abstracts
% =================================================
\newpage
\begin{abstracts}

  \title{Spatially Resolved Aromatic Band Ratios in the Orion Bar with
    JWST/NIRSpec}
  \begin{abstract}[Leiden Observatory, Leiden University \\ Institut
      d'Astrophysique Spatiale]{A.\ Tielens$^1$, M.\ Berne$^2$ and
      K.\ Habart$^2$}
    We present JWST/NIRSpec IFU observations of the Orion Bar photodissociation
    region, mapping the 3.3, 6.2 and 7.7 \textmu m aromatic infrared bands at
    0.2 arcsecond resolution. Band ratios vary smoothly with depth into the
    cloud, tracing the changing PAH ionisation balance predicted by
    photochemical models. We find no evidence for the sharp size-selective
    photodestruction zone proposed in earlier ground-based studies.
  \end{abstract}

  \title{A Survey of PAH Emission in Nearby Star-Forming Dwarf Galaxies}
  \begin{abstract}[Max Planck Institute for Astronomy \\ Princeton
      University]{R.\ Sandstrom$^1$ and L.\ Draine$^2$}
    Using archival Spitzer/IRS spectra of 42 nearby dwarf galaxies, we measure
    the PAH-to-total-infrared luminosity fraction as a function of gas-phase
    metallicity. The fraction declines steeply below 12 + log(O/H) = 8.0,
    consistent with PAH destruction in the harder radiation fields of
    low-metallicity star-forming regions.
  \end{abstract}

\end{abstracts}

% =================================================
%  Meetings
% =================================================
\newpage
\headergraphic{images/Header_Meetings.png}
\section*{~}
\addcontentsline{toc}{section}{Meetings}
\normalsize

\begin{itemize}
  \item \textbf{PAHs and the Universe III} --- 12--16 April 2027,
    Granada, Spain. Abstract deadline: 15 January 2027.
  \item \textbf{Molecular Astrophysics Gordon Research Conference} ---
    June 2027, Newport, RI, USA.
  \item \textbf{JWST Cycle 4 Proposal Deadline} --- 1 November 2026.
\end{itemize}

% =================================================
%  Announcements
% =================================================
\headergraphic{images/Header_Announcements.png}
\section*{~}
\addcontentsline{toc}{section}{Announcements}
\normalsize

\textbf{Postdoctoral position, Leiden Observatory.} A 2-year postdoctoral
position is available to work on laboratory spectroscopy of PAH cations
relevant to JWST mid-infrared observations. Applications due 1 December
2026.

\textbf{New database release.} The PAH IR Spectroscopic Database (v4.0)
now includes computed spectra for over 4000 PAH species, including
nitrogen-substituted and partially dehydrogenated variants.

% =================================================
%  Box at the end
% =================================================
\textcolor{white}{.}
\vfill

\begin{shaded}
\usefont {OT1} {phv}{m}{n}
\fontsize{15}{16}
\selectfont

\centerline{\textbf{\textcolor{white}{AstroPAH Newsletter}}}

\begin{multicols}{2}
\normalsize

\noindent \href{http://astropah-news.strw.leidenuniv.nl}{\textcolor{white}{http://astropah-news.strw.leidenuniv.nl}}

\noindent \href{mailto:astropah@strw.leidenuniv.nl}{\textcolor{white}{astropah@strw.leidenuniv.nl}}

\columnbreak

\flushright \textcolor{white}{Next issue: \nextissue}

\textcolor{white}{Submission deadline: \nextdeadline}

\end{multicols}
\end{shaded}

\end{document}
